\documentclass[11pt,a4paper]{article}
\usepackage{kotex}
\usepackage{amsmath,amssymb,booktabs,geometry}
\geometry{margin=25mm}
\title{I01 선검증 5: 표지 수 요동과 블록 간 표지 분포 변화}
\author{RELATIONAL ATLAS}
\date{2026년 10월 4일}
\begin{document}
\maketitle

\section{판정}
발전연구 주제는 선택하지 않는다. 사전 고정한 Aer 실험은
\emph{무작위 표지 수 요동}과 \emph{표지 분포 자체의 블록 간 변화}를
분리했다. 목표 분포가 같은 두 블록은 표지를 지운 뒤에도
정확한 평균 null이며, 무작위 표지 수는 분산만 늘린다.
반면 과학 위상을 전혀 바꾸지 않아도 균일 분포에서 비균일
분포로 바뀐 비제약 상태는 표지를 지운 관측률이 변했다.
기록된 표지로 두 블록을 같은 균일 목표에 표준화하면 이
차이는 정확 모형에서 0이다. 이는 표지 기록이라는 자원을
사용한 분석이지, 실제로 버린 표지의 정보 회복이 아니다.

\section{사전 설계와 정확 예측}
\texttt{prevalidation\_protocol5\_2026-10-04.tex}에
$n=3$, cutoff 1, $L=7$, $\beta=0.45$, 기준 폭 $\beta_r=0.15$,
$E=E_0+0.37E_0^2/3$, 검출 바이어스 1.25 rad,
$\tau=0,\pi/2$ 및 비제약/$P=0$ 상태를 고정했다.
세 블록의 \emph{과학 공통 위상은 모두 0}이다.
$A,B$의 목표 표지는 균일 $a_t=1/7$,
$C$는 $a_t=[1+0.35\cos(2\pi t/7)]/7$이다.
블록$\times$상태$\times$시계$\times$축마다 $N=229376$개 표지를
독립 다항분포로 뽑았다. 목표는 같아도 $A,B$의 실제 표지 수는 다르다.
배분 seed 20261023, Aer seed 시작값 20261024,
workers 6으로 252회로를 실행했다.

정확 모형의 $A-C$ 비제약 목표-혼합 최대 확률 차이는
$\tau=0,\pi/2$에서 0.02749/0.01343이다.
$P=0$에서는 0이다. $A-B$는 두 목표가 같으므로
무조건부 평균 차이가 0이다. 표지별 성공확률을 $p_t$라 하면
균일 목표의 축별 복귀율 분산은
\begin{align}
\operatorname{Var}_{\rm fixed}(\hat p)
 &=\frac1N\sum_t a_t p_t(1-p_t),\\
\operatorname{Var}_{\rm random}(\hat p)
 &=\frac{\bar p(1-\bar p)}{N}
 =\operatorname{Var}_{\rm fixed}(\hat p)
 +\frac{\operatorname{Var}_{a}(p_t)}{N}.
\end{align}
따라서 표지에 따른 응답 차이가 큰 비제약 상태에서는
무작위 표지 수를 고정 배분으로 오해하면 불확실도를 과소평가한다.
정확 분산 비율은 $\tau=0,\pi/2$에서 1.3549/1.1602,
$P=0$에서는 1이다.

\section{Aer의 두 분석 가족}
첫 가족은 표지를 \emph{삭제}한 성공수만 사용했다.
$A-B,A-C$ 각각에서 두 상태$\times$두 시계를 교차한
8검정을 10000회 조건부 재배분으로 분석했다. 표지는 매 시행
독립 추출됐고 각 축의 총 shots는 같으므로, 목표 분포로 적분한
복귀 결과는 Bernoulli 시행이다. 두 목표-혼합 확률이 같다는
null에서 축별 합친 성공수를 조건화했다. 사전 Bonferroni
문턱은 $0.00625$이다.

\begin{center}
\begin{tabular}{@{}llrr@{}}
\toprule
대조/상태 & $\tau/\pi$ & 표지 삭제 $p$ & 균일 표준화 $p$ \\
\midrule
$A-B$/비제약 & 0 & 0.29027 & 0.59364 \\
$A-B$/비제약 & 1/2 & 0.58164 & 0.41596 \\
$A-B$/$P=0$ & 0 & 0.49835 & 0.51295 \\
$A-B$/$P=0$ & 1/2 & 0.43556 & 0.44456 \\
$A-C$/비제약 & 0 & 0.00010$^{*}$ & 0.89581 \\
$A-C$/비제약 & 1/2 & 0.00010$^{*}$ & 0.04610 \\
$A-C$/$P=0$ & 0 & 0.68473 & 0.84442 \\
$A-C$/$P=0$ & 1/2 & 0.99790 & 0.97090 \\
\bottomrule
\end{tabular}
\end{center}

$^{*}$10000회 재배분에서 관측 통계량을 넘은 표본이 0개여서
$(0+1)/10001$이다. 이는 Monte Carlo 해상도 바닥이지
참 $p$값의 엄밀한 상한이 아니다.
표지 삭제 가족에서는 사전 문턱을 통과한 것이 $A-C$ 비제약
두 시계 조건뿐이었다. $A-B$와 $P=0$ 대조군은 통과하지 않았다.

둘째 가족은 \emph{같은 raw counts}의 표지 기록을 써서
블록별 성공률을 균일 목표 $1/7$로 표준화했다. 표지별 shots
수가 달라도 같은 표지의 두 블록 성공확률은 과학 위상을
바꾸지 않았으므로 같다. 축$\times$표지 층별 총 성공수를
조건화한 정확 hypergeometric 재배분을 사용했다.
이 가족의 8검정은 모두 사전 문턱 0.00625를 넘지 않았다.
다만 $A-C$ 비제약 $\tau=\pi/2$의 명목 $p=0.04610$을
숨기지 않는다. 한 seed의 비검출이 일반적인 보정 성공을
보증하지 않으며, 정확 모형의 표준화 목표가 0이라는 주장과
유한 샷의 우연을 구분한다.

\section{사후 계수 및 오류율 감사}
252개 Aer 회로의 성공수를 각 정확 이항확률로 표준화한 평균은
$-0.1108$, 분산은 $0.9644$였다. 5000개 정확 이항패널의
95\% 평균 범위 $[-0.1213,0.1211]$ 및 분산 범위
$[0.7953,1.1730]$ 안에 있지만 평균은 하단에 가까웠다.
추가 반복 없이 이 한 배치를 이상 또는 편향으로 단정하지 않는다.

사후 기술 감사로 균일 목표의 독립 무작위 표지 배분과 이항
측정을 20000패널씩 생성했다. 비제약/$\tau=0$에서
올바른 무작위 배분 분산을 쓰면 명목 5\% 양성이
983/20000(4.92\%)였지만, 고정 배분 분산을 잘못 쓰면
2488/20000(12.44\%)였다. $\tau=\pi/2$에서는
1048/20000(5.24\%) 대 1641/20000(8.21\%)였다.
$P=0$는 표지별 응답이 같아 두 계산 모두 약 5\%였다.
이 감사는 선언된 토이 분포에서의 오류율 확인이며 다른
장치나 임의의 드리프트 법칙에 대한 보장은 아니다.

\section{적용 범위와 남은 검증}
이번 결과는 목표 표지 분포가 \emph{알려져 있고}, 표지가
독립적으로 뽑히며, 각 축의 배분이 독립이라는 설계에 한정된다.
실제 표지 기록이 없거나 목표 분포 자체가 알려져 있지 않으면
표준화 분석을 그대로 실행할 수 없다. 블록 내 상관된 표지
생성, 기록 오류 및 위상 표지의 보정 오차는 남아 있다.
또 회로는 직접 에너지 위상 유니터리를 사용해 별도 시계
레지스터를 후선택하지 않았다. 다른 cutoff와 실제 시계 회로를
검증하기 전에는 넓은 식별 명제나 발전연구 주제를 정하지 않는다.
QPU도 사용하지 않았다.

\paragraph{재현 자료.}
사전 프로토콜: \texttt{prevalidation\_protocol5\_2026-10-04.tex}.
정확 pilot: \texttt{../src/relational\_atlas\_i01/holdout/qiskit\_rank\_holdout5\_power.py},
\texttt{../results/qiskit\_rank\_holdout5\_power.json}.
Aer raw counts와 두 가족 검정:
\texttt{../src/relational\_atlas\_i01/holdout/qiskit\_rank\_holdout5\_random\_labels.py},
\texttt{../results/qiskit\_rank\_holdout5\_random\_labels.json}.
사후 오류율 감사:
\texttt{../src/relational\_atlas\_i01/holdout/qiskit\_rank\_holdout5\_allocation\_audit.py},
\texttt{../results/qiskit\_rank\_holdout5\_allocation\_audit.json}.
\end{document}
